% 闭卷背诵稿模板。同一份源文件编译两个版本：
%   python scripts/compile_tex.py <文件>.tex          背诵稿：\kw{} 里的关键词加粗
%   python scripts/compile_tex.py <文件>.tex --fill   填空版：\kw{} 变成下划线空白，用来自测
% 规则：只改标题信息和正文，导言区（\begin{document}之前）不要动，不要自己加宏包和宏，不引入外部图片
% 【 】里是要替换的占位内容（替换时连括号一起去掉）。公式外的 % & # _ $ ^ ~ { } \ 要转义，见 references/exam-notes.md
% 写法：背诵优先。按考点权重从高到低分块，一块能在十来分钟背完；\kw{} 只包一个短关键词（不超过8个字），
%       因为填空版的空白不能折行
\documentclass[a4paper,fontset=windows]{ctexart}
\ifdefined\SBSIZE\else\def\SBSIZE{11}\fi
\ifdefined\SBLEAD\else\def\SBLEAD{17}\fi
\ifdefined\SBMARGIN\else\def\SBMARGIN{22}\fi
\usepackage[margin=\SBMARGIN mm]{geometry}
\usepackage{amsmath,amssymb,enumitem,booktabs,xcolor}
\makeatletter
\renewcommand\normalsize{\@setfontsize\normalsize{\SBSIZE}{\SBLEAD}%
  \abovedisplayskip 4\p@ \@plus2\p@ \@minus2\p@
  \belowdisplayskip \abovedisplayskip
  \abovedisplayshortskip 2\p@ \@plus1\p@
  \belowdisplayshortskip 2\p@ \@plus1\p@
  \let\@listi\@listI}
\makeatother
\normalsize
\pagestyle{plain}
\setlength{\parindent}{0pt}
\setlength{\parskip}{4pt}
\setlist{nosep,leftmargin=1.6em,itemsep=3pt}
\ctexset{section={number=\arabic{section},name={,},aftername=\hspace{0.6em},format=\Large\bfseries,beforeskip=14pt,afterskip=6pt},
         subsection={number=\arabic{section}.\arabic{subsection},name={,},aftername=\hspace{0.5em},format=\large\bfseries,beforeskip=8pt,afterskip=3pt}}
\newlength{\sbw}
\newcommand{\kw}[1]{\ifdefined\SBFILL\settowidth{\sbw}{#1}\underline{\hspace*{\dimexpr\sbw+0.6em\relax}}\else\textbf{#1}\fi}

\begin{document}
\begin{center}
{\LARGE\bfseries 【课程名】　闭卷背诵稿\ifdefined\SBFILL（填空版）\fi}\\[4pt]
{\small 考试：【YYYY-MM-DD】　范围：【范围】}
\end{center}

\section{高权重（必背）}
\subsection{【知识点名】}
\begin{enumerate}
\item 【一句话，关键词用 \kw{} 包起来】
\item 【公式：$f(x)$，其中 \kw{【符号含义】}】
\end{enumerate}

% 对比类内容用表格，更好背
\begin{center}
\begin{tabular}{@{}lll@{}}
\toprule
 & 【A】 & 【B】 \\
\midrule
【维度】 & 【\kw{内容}】 & 【\kw{内容}】 \\
\bottomrule
\end{tabular}
\end{center}

\section{中权重}
\subsection{【知识点名】}
\begin{enumerate}
\item 【\kw{关键词}】
\end{enumerate}

\section{低权重（时间够再背）}
\subsection{【知识点名】}
\begin{enumerate}
\item 【\kw{关键词}】
\end{enumerate}

\end{document}
